\begin{pastexam}[label={px:ntu-111-dsaa-04}]{臺大資工所 111 學年度 第 4 題}
（單選。 第 1--5 題共用下面的說明和表格）

For each of the following algorithms in the next five problems, what is the tightest asymptotic
upper bound for its runtime complexity for $n$ numbers? Please use the following table of (A)-(E)
choices when answering.
\begin{center}
\begin{tabular}{c|c|c|c|c}
(A) & (B) & (C) & (D) & (E) \\ \hline
$O(1)$ & $O(\lg n)$ & $O(n)$ & $O(n\lg n)$ & $O(n^2)$
\end{tabular}
\end{center}

(3 points) Quick sort: worst-case time?
\end{pastexam}

\begin{pastsoln}
答案是 (E) $O(n^2)$。

快速排序（資結書第~14.5~節）每次分割花 $\Theta(m)$（$m$ 是這一段的長度）。 最壞的情況是每次
基準都是最大或最小值，切成 $m - 1$ 和 $0$ 個（例如取最後一個當基準、輸入已經排好）：
\[
T(n) = T(n - 1) + \Theta(n) = \Theta\Bigl(\sum_{m=1}^{n} m\Bigr) = \Theta(n^2).
\]
這種輸入確實存在，所以最壞情況的最緊上界是 $O(n^2)$。

(D) $O(n \lg n)$ 是平均（或隨機選基準時的期望）時間，不是最壞情況。 (A)--(C) 更小。
\end{pastsoln}
